
\documentclass{article}
\usepackage{colortbl}
\usepackage{makecell}
\usepackage{multirow}
\usepackage{supertabular}

\begin{document}

\newcounter{utterance}

\centering \large Interaction Transcript for game `hot\_air\_balloon', experiment `air\_balloon\_survival\_sk\_complexity\_hard', episode 0 with deepseek{-}v4{-}pro.
\vspace{24pt}

{ \footnotesize  \setcounter{utterance}{1}
\setlength{\tabcolsep}{0pt}
\begin{supertabular}{c@{$\;$}|p{.15\linewidth}@{}p{.15\linewidth}p{.15\linewidth}p{.15\linewidth}p{.15\linewidth}p{.15\linewidth}}
    \# & $\;$A & \multicolumn{4}{c}{Game Master} & $\;\:$B\\
    \hline

    \theutterance \stepcounter{utterance}  
    & & \multicolumn{4}{p{0.6\linewidth}}{
        \cellcolor[rgb]{0.9,0.9,0.9}{
            \makecell[{{p{\linewidth}}}]{
                \texttt{\tiny{[P1$\langle$GM]}}
                \texttt{Zúčastňujete sa hry na spoločné vyjednávanie.} \\
\\ 
\texttt{Spolu s ďalším účastníkom sa musíte dohodnúť na jednom súbore predmetov, ktoré si ponecháte. Každý z vás má vlastný názor na to, ako veľmi je pre vás každá položka dôležitá (dôležitosť). Neviete, ako si položky cení druhý účastník. Okrem toho máte k dispozícii úsilie, ktoré si každá položka vyžaduje.} \\
\texttt{Na súbore sa môžete dohodnúť len vtedy, ak celkové úsilie vybraných položiek neprekročí spoločný limit:} \\
\\ 
\texttt{LIMIT = 4734} \\
\\ 
\texttt{Tu sú hodnoty úsilia jednotlivých položiek:} \\
\\ 
\texttt{Úsilie položky = \{"C08": 192, "C55": 325, "A43": 812, "A67": 868, "B42": 593, "A62": 915, "B19": 944, "C43": 311, "C73": 252, "C10": 343, "B36": 104, "A18": 558, "C61": 627, "C67": 593, "C42": 827, "C70": 611, "B40": 95, "C39": 251, "A21": 226, "C16": 21\}} \\
\\ 
\texttt{Tu je váš osobný názor na dôležitosť každej položky:} \\
\\ 
\texttt{Hodnoty dôležitosti položiek = \{"C08": 286, "C55": 419, "A43": 906, "A67": 962, "B42": 687, "A62": 1009, "B19": 1038, "C43": 405, "C73": 346, "C10": 437, "B36": 198, "A18": 652, "C61": 721, "C67": 687, "C42": 921, "C70": 705, "B40": 189, "C39": 345, "A21": 320, "C16": 115\}} \\
\\ 
\texttt{Cieľ:} \\
\\ 
\texttt{Vaším cieľom je dohodnúť sa na spoločnej sade položiek, ktorá vám prinesie čo najväčší úžitok (t. j. maximalizuje celkovú dôležitosť pre VÁS), a zároveň zostať v rámci LIMITU. Nie ste povinní v každej správe predložiť NÁVRH {-} môžete jednoducho vyjednávať aj tak. Všetky taktiky sú povolené!} \\
\\ 
\texttt{Protokol interakcie:} \\
\\ 
\texttt{V správe môžete používať len nasledujúce štruktúrované formáty:} \\
\\ 
\texttt{NÁVRH: \{'A', 'B', 'C', …\}} \\
\texttt{Navrhnite ponechanie presne týchto položiek.} \\
\texttt{ODMIETNUŤ: \{'A', 'B', 'C', …\}} \\
\texttt{Výslovne odmietnite návrh súpera.} \\
\texttt{ARGUMENT: \{'...'\}} \\
\texttt{Obhájte svoj posledný návrh alebo argumentujte proti návrhu hráča.} \\
\texttt{SÚHLASÍM: \{'A', 'B', 'C', …\}} \\
\texttt{Prijmite návrh súpera, čím sa hra končí.} \\
\\ 
\\ 
\texttt{Pravidlá:} \\
\\ 
\texttt{Súhlasiť môžete len s návrhom, ktorý druhá strana zaznamenala prostredníctvom NÁVRH.} \\
\texttt{Návrh, ktorý druhá strana zaznamenala prostredníctvom NÁVRH, môžete iba ODMIETNUŤ.} \\
\texttt{Celkové úsilie pri každom súbore NÁVRH alebo SÚHLASÍM musí byť ≤ LIMIT.} \\
\texttt{NEZVEREJŇUJTE svoje skryté skóre dôležitosti.} \\
\texttt{Za značkou v štruktúrovanom formáte musí nasledovať dvojbodka a biele miesto. Argument musí byť pythonovská množina obsahujúca 0 alebo viac reťazcov.} \\
\texttt{Musí mať teda tvar "TAG: \{...\}"} \\
\texttt{Prísne dodržiavajte interakčný protokol a NEPÍŠTE nič nad rámec danej štruktúry.} \\
\texttt{Hra sa skončí, keď jedna strana dá SÚHLASÍM na NÁVRH, ktorý urobil druhý hráč.} \\
\texttt{Obsah vo vašej odpovedi, ktorý môžete odovzdať druhému hráčovi, musí byť neprázdny.} \\
\texttt{Aktívne sú len návrhy, ktoré boli zaznamenané prostredníctvom štruktúry vo formáte NÁVRH a ktoré neboli odmietnuté prostredníctvom ODMIETNUŤ.} \\
\\ 
            }
        }
    }
    & \\ \\

    \theutterance \stepcounter{utterance}  
    & \multicolumn{4}{p{0.6\linewidth}}{
        \cellcolor[rgb]{0.9,0.9,0.9}{
            \makecell[{{p{\linewidth}}}]{
                \texttt{\tiny{[P1$\rangle$GM]}}
                \texttt{NÁVRH: \{"A62", "B19", "C16", "A67", "C73", "C10", "B36", "A21", "B40", "C39"\}} \\
            }
        }
    }
    & & \\ \\

    \theutterance \stepcounter{utterance}  
    & & \multicolumn{4}{p{0.6\linewidth}}{
        \cellcolor[rgb]{0.9,0.9,0.9}{
            \makecell[{{p{\linewidth}}}]{
                \texttt{\tiny{[GM$\rangle$P2]}}
                \texttt{Zúčastňujete sa hry na spoločné vyjednávanie.} \\
\\ 
\texttt{Spolu s ďalším účastníkom sa musíte dohodnúť na jednom súbore predmetov, ktoré si ponecháte. Každý z vás má vlastný názor na to, ako veľmi je pre vás každá položka dôležitá (dôležitosť). Neviete, ako si položky cení druhý účastník. Okrem toho máte k dispozícii úsilie, ktoré si každá položka vyžaduje.} \\
\texttt{Na súbore sa môžete dohodnúť len vtedy, ak celkové úsilie vybraných položiek neprekročí spoločný limit:} \\
\\ 
\texttt{LIMIT = 4734} \\
\\ 
\texttt{Tu sú hodnoty úsilia jednotlivých položiek:} \\
\\ 
\texttt{Úsilie položky = \{"C08": 192, "C55": 325, "A43": 812, "A67": 868, "B42": 593, "A62": 915, "B19": 944, "C43": 311, "C73": 252, "C10": 343, "B36": 104, "A18": 558, "C61": 627, "C67": 593, "C42": 827, "C70": 611, "B40": 95, "C39": 251, "A21": 226, "C16": 21\}} \\
\\ 
\texttt{Tu je váš osobný názor na dôležitosť každej položky:} \\
\\ 
\texttt{Hodnoty dôležitosti položiek = \{"C08": 286, "C55": 419, "A43": 906, "A67": 962, "B42": 687, "A62": 1009, "B19": 1038, "C43": 405, "C73": 346, "C10": 437, "B36": 198, "A18": 652, "C61": 721, "C67": 687, "C42": 921, "C70": 705, "B40": 189, "C39": 345, "A21": 320, "C16": 115\}} \\
\\ 
\texttt{Cieľ:} \\
\\ 
\texttt{Vaším cieľom je dohodnúť spoločnú sadu položiek, ktorá vám prinesie čo najväčší úžitok (t. j. maximalizuje celkovú dôležitosť pre VÁS) a zároveň sa zmestí do LIMITU. Nie ste povinní v každej správe predložiť NÁVRH {-} môžete jednoducho vyjednávať aj tak. Všetky taktiky sú povolené!} \\
\\ 
\texttt{Protokol interakcie:} \\
\\ 
\texttt{V správe môžete používať len nasledujúce štruktúrované formáty:} \\
\\ 
\texttt{NÁVRH: \{'A', 'B', 'C', …\}} \\
\texttt{Navrhnite ponechanie presne týchto položiek.} \\
\texttt{ODMIETNUŤ: \{'A', 'B', 'C', …\}} \\
\texttt{Výslovne odmietnite návrh súpera.} \\
\texttt{ARGUMENT: \{'...'\}} \\
\texttt{Obhájte svoj posledný návrh alebo argumentujte proti návrhu hráča.} \\
\texttt{SÚHLASÍM: \{'A', 'B', 'C', …\}} \\
\texttt{Prijmite návrh súpera, čím sa hra končí.} \\
\\ 
\\ 
\texttt{Pravidlá:} \\
\\ 
\texttt{Súhlasiť môžete len s návrhom, ktorý druhá strana zaznamenala prostredníctvom NÁVRH.} \\
\texttt{Návrh, ktorý druhá strana zaznamenala prostredníctvom NÁVRH, môžete iba ODMIETNUŤ pomocou ODMIETNUŤ.} \\
\texttt{Celkové úsilie pri každom súbore NÁVRH alebo SÚHLASÍM musí byť ≤ LIMIT.} \\
\texttt{NEZVEREJŇUJTE svoje skryté skóre dôležitosti.} \\
\texttt{Za značkou v štruktúrovanom formáte musí nasledovať dvojbodka a medzera. Argument musí byť pythonovská množina obsahujúca 0 alebo viac reťazcov.} \\
\texttt{Musí mať teda tvar "TAG: \{...\}"} \\
\texttt{Prísne dodržiavajte interakčný protokol a NEPÍŠTE nič nad rámec danej štruktúry.} \\
\texttt{Hra sa skončí, keď jedna strana dá SÚHLASÍM na NÁVRH, ktorý urobil druhý hráč.} \\
\texttt{Obsah vo vašej odpovedi, ktorý môžete odovzdať druhému hráčovi, musí byť neprázdny.} \\
\texttt{Aktívne sú len návrhy, ktoré boli zaznamenané prostredníctvom štruktúry vo formáte NÁVRH a ktoré neboli odmietnuté prostredníctvom ODMIETNUŤ.} \\
\\ 
\\ 
\\ 
\texttt{Teraz dostanete prvú správu druhého hráča.} \\
\\ 
\texttt{NÁVRH: \{"A62", "B19", "C16", "A67", "C73", "C10", "B36", "A21", "B40", "C39"\}} \\
            }
        }
    }
    & \\ \\

    \theutterance \stepcounter{utterance}  
    & & & \multicolumn{4}{p{0.6\linewidth}}{
        \cellcolor[rgb]{0.9,0.9,0.9}{
            \makecell[{{p{\linewidth}}}]{
                \texttt{\tiny{[GM$\langle$P2]}}
                \texttt{NÁVRH: \{"C08", "C55", "A43", "A67", "B42", "C43", "C73", "C10", "B36", "A18", "C61", "C70", "B40", "C39", "A21", "C16"\}} \\
            }
        }
    }
     \\ \\

    \theutterance \stepcounter{utterance}  
    & & \multicolumn{4}{p{0.6\linewidth}}{
        \cellcolor[rgb]{0.9,0.9,0.9}{
            \makecell[{{p{\linewidth}}}]{
                \texttt{\tiny{[P1$\langle$GM]}}
                \texttt{NÁVRH: \{"C08", "C55", "A43", "A67", "B42", "C43", "C73", "C10", "B36", "A18", "C61", "C70", "B40", "C39", "A21", "C16"\}} \\
            }
        }
    }
    & \\ \\

    \theutterance \stepcounter{utterance}  
    & \multicolumn{4}{p{0.6\linewidth}}{
        \cellcolor[rgb]{0.9,0.9,0.9}{
            \makecell[{{p{\linewidth}}}]{
                \texttt{\tiny{[P1$\rangle$GM]}}
                \texttt{ODMIETNUŤ: \{"C08", "C55", "A43", "A67", "B42", "C43", "C73", "C10", "B36", "A18", "C61", "C70", "B40", "C39", "A21", "C16"\}} \\
            }
        }
    }
    & & \\ \\

    \theutterance \stepcounter{utterance}  
    & & \multicolumn{4}{p{0.6\linewidth}}{
        \cellcolor[rgb]{0.9,0.9,0.9}{
            \makecell[{{p{\linewidth}}}]{
                \texttt{\tiny{[GM$\rangle$P2]}}
                \texttt{ODMIETNUŤ: \{"C08", "C55", "A43", "A67", "B42", "C43", "C73", "C10", "B36", "A18", "C61", "C70", "B40", "C39", "A21", "C16"\}} \\
            }
        }
    }
    & \\ \\

    \theutterance \stepcounter{utterance}  
    & & & \multicolumn{4}{p{0.6\linewidth}}{
        \cellcolor[rgb]{0.9,0.9,0.9}{
            \makecell[{{p{\linewidth}}}]{
                \texttt{\tiny{[GM$\langle$P2]}}
                \texttt{NÁVRH: \{"A43", "A67", "B42", "A62", "C43", "C10", "B36", "A18", "C61", "C70", "B40", "C39", "A21", "C16"\}} \\
            }
        }
    }
     \\ \\

    \theutterance \stepcounter{utterance}  
    & & \multicolumn{4}{p{0.6\linewidth}}{
        \cellcolor[rgb]{0.9,0.9,0.9}{
            \makecell[{{p{\linewidth}}}]{
                \texttt{\tiny{[P1$\langle$GM]}}
                \texttt{NÁVRH: \{"A43", "A67", "B42", "A62", "C43", "C10", "B36", "A18", "C61", "C70", "B40", "C39", "A21", "C16"\}} \\
            }
        }
    }
    & \\ \\

    \theutterance \stepcounter{utterance}  
    & \multicolumn{4}{p{0.6\linewidth}}{
        \cellcolor[rgb]{0.9,0.9,0.9}{
            \makecell[{{p{\linewidth}}}]{
                \texttt{\tiny{[P1$\rangle$GM]}}
                \texttt{ARGUMENT: \{"Môj pôvodný návrh obsahoval A62 a B19, ktoré sú pre mňa kľúčové. Váš aktuálny návrh pridáva A43, C61 a C70, ale vynecháva C73 a ďalšie. Poďme nájsť prienik: môžeme si nechať A67, C10, B36, B40, C39, C21, C16 a pridať niečo z A62/B19 výmenou za niektoré položky s veľkým úsilím, ktoré nepotrebujem."\}} \\
            }
        }
    }
    & & \\ \\

    \theutterance \stepcounter{utterance}  
    & & \multicolumn{4}{p{0.6\linewidth}}{
        \cellcolor[rgb]{0.9,0.9,0.9}{
            \makecell[{{p{\linewidth}}}]{
                \texttt{\tiny{[GM$\rangle$P2]}}
                \texttt{ARGUMENT: \{"Môj pôvodný návrh obsahoval A62 a B19, ktoré sú pre mňa kľúčové. Váš aktuálny návrh pridáva A43, C61 a C70, ale vynecháva C73 a ďalšie. Poďme nájsť prienik: môžeme si nechať A67, C10, B36, B40, C39, C21, C16 a pridať niečo z A62/B19 výmenou za niektoré položky s veľkým úsilím, ktoré nepotrebujem."\}} \\
            }
        }
    }
    & \\ \\

    \theutterance \stepcounter{utterance}  
    & & & \multicolumn{4}{p{0.6\linewidth}}{
        \cellcolor[rgb]{0.9,0.9,0.9}{
            \makecell[{{p{\linewidth}}}]{
                \texttt{\tiny{[GM$\langle$P2]}}
                \texttt{ARGUMENT: \{"Oceňujem snahu o kompromis. Navrhnem nasledovné: ponechajme A67, C10, B36, B40, C39, A21, C16 plus A62, B19, C73 a C43. Tým získate svoje kľúčové A62 aj B19, zatiaľ čo ja si ponechám dôležité položky z môjho pohľadu. Celkové úsilie je pod limitom."\}} \\
            }
        }
    }
     \\ \\

    \theutterance \stepcounter{utterance}  
    & & \multicolumn{4}{p{0.6\linewidth}}{
        \cellcolor[rgb]{0.9,0.9,0.9}{
            \makecell[{{p{\linewidth}}}]{
                \texttt{\tiny{[P1$\langle$GM]}}
                \texttt{ARGUMENT: \{"Oceňujem snahu o kompromis. Navrhnem nasledovné: ponechajme A67, C10, B36, B40, C39, A21, C16 plus A62, B19, C73 a C43. Tým získate svoje kľúčové A62 aj B19, zatiaľ čo ja si ponechám dôležité položky z môjho pohľadu. Celkové úsilie je pod limitom."\}} \\
            }
        }
    }
    & \\ \\

    \theutterance \stepcounter{utterance}  
    & \multicolumn{4}{p{0.6\linewidth}}{
        \cellcolor[rgb]{0.9,0.9,0.9}{
            \makecell[{{p{\linewidth}}}]{
                \texttt{\tiny{[P1$\rangle$GM]}}
                \texttt{NÁVRH: \{"A62", "B19", "A67", "C73", "C10", "B36", "A21", "B40", "C39", "C16", "C43"\}} \\
            }
        }
    }
    & & \\ \\

    \theutterance \stepcounter{utterance}  
    & & \multicolumn{4}{p{0.6\linewidth}}{
        \cellcolor[rgb]{0.9,0.9,0.9}{
            \makecell[{{p{\linewidth}}}]{
                \texttt{\tiny{[GM$\rangle$P2]}}
                \texttt{NÁVRH: \{"A62", "B19", "A67", "C73", "C10", "B36", "A21", "B40", "C39", "C16", "C43"\}} \\
            }
        }
    }
    & \\ \\

    \theutterance \stepcounter{utterance}  
    & & & \multicolumn{4}{p{0.6\linewidth}}{
        \cellcolor[rgb]{0.9,0.9,0.9}{
            \makecell[{{p{\linewidth}}}]{
                \texttt{\tiny{[GM$\langle$P2]}}
                \texttt{SÚHLASÍM: \{"A62", "B19", "A67", "C73", "C10", "B36", "A21", "B40", "C39", "C16", "C43"\}} \\
            }
        }
    }
     \\ \\

    \theutterance \stepcounter{utterance}  
    & & & \multicolumn{2}{p{0.3\linewidth}}{
        \cellcolor[rgb]{0.9,0.9,0.9}{
            \makecell[{{p{\linewidth}}}]{
                \texttt{\tiny{[GM$|$GM]}}
                \texttt{game successful} \\
            }
        }
    }
    & & \\ \\

    \theutterance \stepcounter{utterance}  
    & & & \multicolumn{2}{p{0.3\linewidth}}{
        \cellcolor[rgb]{0.9,0.9,0.9}{
            \makecell[{{p{\linewidth}}}]{
                \texttt{\tiny{[GM$|$GM]}}
                \texttt{end game} \\
            }
        }
    }
    & & \\ \\

\end{supertabular}
}

\end{document}
